% !TeX root = BookN.tex


\title{Elastic Analysis of Finite Domains with Corner Points and Eigenstrains}

\titlerunning{Finite Domains with Corner Points and Eigenstrains}

\author{Mykola Yuzvyak and Yuriy Tokovyy}

\institute{
	Mykola Yuzvyak \at 
	\AffPidstryhach \\
	\email{yuzvyak@iapmm.lviv.ua}
	\and
	Yuriy Tokovyy \at 
	\AffPidstryhach \\
	\email{tokovyy@iapmm.lviv.ua}
}

\maketitle

\abstract{
	A generalized {direct-integration method} is presented for two-dimensional
	elasticity and thermoelasticity problems in a rectangular domain whose
	elastic moduli are arbitrary functions of both in-plane coordinates.
	The generalization hinges on two ideas.  
	First, all stress components are expressed through a single scalar
	potential—the {Vihak function}.  
	Second, the {method of additional strains} transfers the spatial
	variation of the material constants into initially unknown eigen-strain
	fields, leaving the governing operator with constant coefficients.
	The resulting boundary-value problem for the Vihak function reduces to a
	sequence of second-order integro-differential equations
	supplemented by integral side conditions that follow systematically from
	the original boundary data on the stresses.
	This framework yields closed-form representations for the stresses and
	is readily implemented in an iterative scheme for strongly
	inhomogeneous or layered materials.}%


\graphicspath{{38_Yuzvyak/}}

\section{Introduction}

The fracture behaviour of structural materials is governed by an interplay of physical, mechanical, and geometric factors \cite{Yuzvyak20}.  
Geometric factors are primarily associated with surface stress concentrators such as corner lines and points.  
Within a mathematical model these concentrators coincide with the intersection lines or points of dissimilar coordinate surfaces \cite{Yuzvyak3}.  
The stress field in their neighbourhood differs markedly from that in the surrounding volume, and its evaluation is complicated by discontinuous boundary conditions and surface singularities \cite{Yuzvyak10}.

Spatially varying material properties further influence the mechanical response and fracture resistance of components \cite{Yuzvyak8,Yuzvyak15,Yuzvyak18}.  
Inhomogeneous solids—whose elastic moduli change continuously or step-wise within the body—are increasingly employed in aerospace, biomedical, and structural applications, where tailored gradients enable an improved compromise between competing performance requirements \cite{Yuzvyak9,Yuzvyak27}.  
Comprehensive surveys of fabrication techniques, application areas, and modelling strategies for functionally graded materials (FGMs) are provided in \cite{Yuzvyak6,Yuzvyak13}, while the most recent review dedicated to additively manufactured FGMs emphasises optimisation of gradient profiles for internal-stress control \cite{Yuzvyak1}.

Exact analytical solutions of linear elasticity for non-homogeneous media are scarce because the governing equations contain position-dependent coefficients composed of elastic-moduli profiles \cite{Yuzvyak27,Yuzvyak28}.  
Since these coefficients do not commute with differential operators, many classical techniques become inapplicable \cite{Yuzvyak19}.  
Closed-form results have therefore been obtained only for a few special gradients and geometries.  
Considerable attention has been devoted to media with one-dimensional property variation—typical of FGMs and multilayer assemblies—where separation of variables remains possible \cite{Yuzvyak14,Yuzvyak25}.  
Problems involving oblique \cite{Yuzvyak11,Yuzvyak12,Yuzvyak17} or fully two-/three-dimensional gradients \cite{Yuzvyak34} are substantially more demanding.

This chapter presents an algorithm for the plane-strain analysis of a rectangular domain whose elastic moduli are arbitrary functions of both in-plane coordinates.  
The problem combines two major challenges: corner singularities and multi-directional inhomogeneity.  
The procedure integrates the direct integration method \cite{Yuzvyak35} with the method of additional (eigen)strains, also referred to as the flexibility method \cite{Yuzvyak2,Yuzvyak5,Yuzvyak21,Yuzvyak22}.  
Position-dependent stiffness terms in the constitutive relations are treated as unknown additional strains, whereas the basic elastic strains are evaluated with constant material parameters.  
An explicit solution is first derived for the homogeneous auxiliary problem with prescribed additional strains by means of the direct integration method, which is particularly suited to domains containing corner points.  
The additional strains are then determined iteratively through a successive-approximation scheme, each step requiring the solution of a system with constant coefficients.

\section{Governing Equations of the Problem}

\noindent
A plane problem of the elasticity theory is considered for the rectangle  
\(
\mathcal{D}=\{(x,y)\in[-a_x,a_x]\times[-a_y,a_y]\},
\)
whose elastic moduli are prescribed as arbitrary functions of the in-plane coordinates \(x\) and \(y\).
The Cartesian stresses \(\sigma_{xx},\sigma_{yy},\sigma_{xy}\) obey the equilibrium relations
\begin{equation}
	\frac{\partial\sigma_{xx}}{\partial x}+\frac{\partial\sigma_{xy}}{\partial y}=0, 
	\qquad
	\frac{\partial\sigma_{xy}}{\partial x}+\frac{\partial\sigma_{yy}}{\partial y}=0,
	\label{YuzvyakEq1}
\end{equation}
together with eight boundary conditions on \(\partial\mathcal{D}\),
\begin{equation}
	\begin{array}{cc}
		\sigma_{xx}(\pm a_x,y)=-p_x^{\pm}(y), & \quad
		\sigma_{xy}(\pm a_x,y)=\tau_{xy}^{\pm}(y),\\[2pt]
		\sigma_{yy}(x,\pm a_y)=-p_y^{\pm}(x), & \quad
		\sigma_{xy}(x,\pm a_y)=\tau_{yx}^{\pm}(x),
	\end{array}
	\label{YuzvyakEq2}
\end{equation}
where \(p_x^{\pm},p_y^{\pm}\) and \(\tau_{xy}^{\pm},\tau_{yx}^{\pm}\) are prescribed normal and tangential tractions.

The total strains comprise elastic and thermal parts.  
With Young’s modulus \(E(x,y)\), Poisson’s ratio \(\nu(x,y)\), and thermal-expansion coefficient \(\alpha_T(x,y)\), the constitutive relations for the \emph{plane-stress} hypothesis read
\begin{equation}
	\varepsilon_{xx}= \frac{\sigma_{xx}-\nu\sigma_{yy}}{E}+\varepsilon_T,\quad
	\varepsilon_{yy}= \frac{\sigma_{yy}-\nu\sigma_{xx}}{E}+\varepsilon_T,\quad
	\varepsilon_{xy}= \frac{1+\nu}{E} \sigma_{xy},
	\label{YuzvyakEq3}
\end{equation}
with \(\varepsilon_T=\alpha_T T(x,y)\) produced by the stationary (or quasi-stationary) temperature field \(T(x,y)\).
For the \emph{plane-strain} formulation one substitutes
\(
E \to  E/(1-\nu^{2}),\;
\nu \to \nu/(1+\nu),\;
\varepsilon_T \to (1+\nu) \varepsilon_T.
\)

Compatibility of strains imposes
\begin{equation}
	2 \frac{\partial^{2}\varepsilon_{xy}}{\partial x \partial y}=
	\frac{\partial^{2}\varepsilon_{xx}}{\partial y^{2}}+
	\frac{\partial^{2}\varepsilon_{yy}}{\partial x^{2}},
	\label{YuzvyakEq4}
\end{equation}
and substitution of Eq.~\eqref{YuzvyakEq3} into \eqref{YuzvyakEq4}, followed by use of Eq.~\eqref{YuzvyakEq1}, yields a compatibility equation expressed through stresses.  
For inhomogeneous media this equation is nonlinear because derivatives of the spatially varying moduli appear explicitly, rendering direct analytical treatment difficult.

To bypass this difficulty the \emph{method of additional (eigen)strains} is employed.  
Let
 \[E_0=\displaystyle\max_{(x,y)\in\mathcal{D}}E(x,y)\]
 and let \(\nu_{0}\) denote the value of \(\nu\) taken at the point where \(E(x,y)\) attains its maximum.
   
The constitutive relations are rewritten as
\begin{equation}
	\begin{array}{ll}
		\varepsilon_{xx}&=\dfrac{\sigma_{xx}-\nu_0\sigma_{yy}}{E_0}+\varepsilon_{xx}^{0},\quad
		\varepsilon_{yy}=\dfrac{\sigma_{yy}-\nu_0\sigma_{xx}}{E_0}+\varepsilon_{yy}^{0},
		\\
		\varepsilon_{xy}&=\dfrac{1+\nu_0}{E_0} \sigma_{xy}+\varepsilon_{xy}^{0},
	\label{YuzvyakEq5}
	\end{array}
\end{equation}
where the additional strains  
\begin{align}
	\varepsilon_{xx}^{0}&=\varepsilon_T+\left(\frac{1}{E}-\frac{1}{E_0}\right)\sigma_{xx}-\left(\frac{\nu}{E}-\tfrac{\nu_0}{E_0}\right)\sigma_{yy},\nonumber\\
	\varepsilon_{yy}^{0}&=\varepsilon_T+\left(\frac{1}{E}-\frac{1}{E_0}\right)\sigma_{yy}-\left(\frac{\nu}{E}-\frac{\nu_0}{E_0}\right)\sigma_{xx},\label{YuzvyakEq6}
	\\
	\varepsilon_{xy}^{0}&=\left(\frac{1+\nu}{E}-\frac{1+\nu_0}{E_0}\right))\sigma_{xy},\nonumber
\end{align}
collect all effects of material inhomogeneity and temperature.

Eliminating strains from Eqs.~\eqref{YuzvyakEq1}, \eqref{YuzvyakEq4}, and \eqref{YuzvyakEq5} gives the governing equation
\begin{equation}
	\Delta\left(\sigma_{xx}+\sigma_{yy}\right)=
	E_0 \left(
	2\frac{\partial^{2}\varepsilon_{xy}^{0}}{\partial x\partial y}
	-\frac{\partial^{2}\varepsilon_{xx}^{0}}{\partial y^{2}}
	-\frac{\partial^{2}\varepsilon_{yy}^{0}}{\partial x^{2}}
	\right),
	\quad
	\Delta=\frac{\partial^{2}}{\partial x^{2}}+\frac{\partial^{2}}{\partial y^{2}},
	\label{YuzvyakEq7}
\end{equation}
which is linear in the stresses but contains the yet unknown additional strains on the right-hand side.


Equation \eqref{YuzvyakEq7} is solved iteratively.  
At every approximation the \(\varepsilon_{ij}{}^{0}\) are regarded as independent functions; the corresponding constant-coefficient problem is tackled analytically through the direct-integration method with the aid of Vihak’s functions, whose use is particularly convenient in rectangles with corner points.  
The updated stresses are then substituted into Eq.~\eqref{YuzvyakEq6} to refine \(\varepsilon_{ij}{}^{0}\); the cycle is repeated until convergence.  
In this way both the influence of corners and the nonlinear impact of arbitrary spatial variation of the elastic moduli are captured within a fully analytic framework.

\section{Vihak Function}

When the compatibility relation \eqref{YuzvyakEq7} was obtained, the
equilibrium equations \eqref{YuzvyakEq1} were differentiated to remove
the shear stress.  On a finite domain this operation is admissible only
if the resulting derivatives satisfy additional boundary constraints
\cite{Yuzvyak16,Yuzvyak26}.  For the rectangle
\(\mathcal{D}=[-a_x,a_x] \times [-a_y,a_y]\) these read  
\begin{equation}
	\left.\frac{\partial\sigma_{xx}}{\partial x}\right|_{x=\pm a_x}= -\frac{d\tau_{xy}^{\pm}(y)}{dy},\qquad
	\left.\frac{\partial\sigma_{yy}}{\partial y}\right|_{y=\pm a_y}= -\frac{d\tau_{yx}^{\pm}(x)}{dx},
	\label{YuzvyakEq8}
\end{equation}
and are equivalent to the traction conditions
\((\ref{YuzvyakEq2})_2\) and \((\ref{YuzvyakEq2})_4\).

Differentiating \((\ref{YuzvyakEq1})_1\) with respect to \(x\) and
\((\ref{YuzvyakEq1})_2\) with respect to \(y\) yields
\[
\frac{\partial^{2}\sigma_{xx}}{\partial x^{2}}=-\frac{\partial^{2}\sigma_{xy}}{\partial x \partial y},\qquad
\frac{\partial^{2}\sigma_{yy}}{\partial y^{2}}=-\frac{\partial^{2}\sigma_{xy}}{\partial x \partial y}.
\]
Hence we define the \emph{Vihak function}
\begin{equation}
	\mathcal{V}(x,y)=
	\frac{\partial^{2}\sigma_{xx}}{\partial x^{2}}
	=\frac{\partial^{2}\sigma_{yy}}{\partial y^{2}}
	=-\frac{\partial^{2}\sigma_{xy}}{\partial x \partial y},
	\label{YuzvyakEq9}
\end{equation}
after \cite{Yuzvyak29,Yuzvyak35}.  Once \(\mathcal{V}\) is fixed, the
stress components satisfy the Poisson problems  
\begin{align}
	\frac{\partial^{2}\sigma_{xy}}{\partial x \partial y}&= -\mathcal{V},
	&\text{with BCs }(\ref{YuzvyakEq2})_{2,4},
	\label{YuzvyakEq10}
	\\
	\frac{\partial^{2}\sigma_{xx}}{\partial x^{2}}&=\;\;\mathcal{V},
	&\text{with BCs }(\ref{YuzvyakEq2})_1,\;(\ref{YuzvyakEq8})_1,
	\label{YuzvyakEq11}
	\\
	\frac{\partial^{2}\sigma_{yy}}{\partial y^{2}}&=\;\;\mathcal{V},
	&\text{with BCs }(\ref{YuzvyakEq2})_3,\;(\ref{YuzvyakEq8})_2.
	\label{YuzvyakEq12}
\end{align}
Each is a two-dimensional Poisson equation with mixed boundary data; the
next subsections integrate these relations to express the full stress
field in terms of~\(\mathcal{V}\).

%-----------------------------------------------------------------
%--- Shear-stress integration and sectional constraints (calculus notation)
Integrating Eq.~\eqref{YuzvyakEq10} with respect to \(x\) and using the
boundary data \((\ref{YuzvyakEq2})_2\) gives
\begin{equation}
	2\frac{\partial \sigma_{xy}(x,y)}{\partial y}=
	\frac{d\tau_{xy}^{+}(y)}{dy}+\frac{d\tau_{xy}^{-}(y)}{dy}
	-\int_{-a_x}^{a_x}\mathcal{V}(\xi,y) \operatorname{sgn}(x-\xi) d\xi ,
	\label{YuzvyakEq13}
\end{equation}
which is compatible with \((\ref{YuzvyakEq2})_2\) provided
\begin{equation}
	\int_{-a_x}^{a_x}\mathcal{V}(x,y) dx=
	\frac{d\tau_{xy}^{-}(y)}{dy}-\frac{d\tau_{xy}^{+}(y)}{dy}.
	\label{YuzvyakEq14}
\end{equation}

A single integration of Eq.~\eqref{YuzvyakEq10} with respect to \(y\) and
the boundary data \((\ref{YuzvyakEq2})_4\) yields
\begin{equation}
	2\frac{\partial \sigma_{xy}(x,y)}{\partial x}=
	\frac{d\tau_{yx}^{+}(x)}{dx}+\frac{d\tau_{yx}^{-}(x)}{dx}
	-\int_{-a_y}^{a_y}\mathcal{V}(x,\eta) \operatorname{sgn}(y-\eta) d\eta ,
	\label{YuzvyakEq15}
\end{equation}
and hence
\begin{equation}
	\int_{-a_y}^{a_y}\mathcal{V}(x,y) dy=
	\frac{d\tau_{yx}^{-}(x)}{dx}-\frac{d\tau_{yx}^{+}(x)}{dx}.
	\label{YuzvyakEq16}
\end{equation}

%-----------------------------------------------------------------
%--- Domain resultant and corner coupling
Integrating Eq.~\eqref{YuzvyakEq14} over \(y\) and Eq.~\eqref{YuzvyakEq16}
over \(x\) gives
\begin{align}
	\iint_{\mathcal{D}}\mathcal{V}(x,y) dx dy&=
	\tau_{yx}^{+}(-a_x)-\tau_{yx}^{+}(a_x)+\tau_{yx}^{-}(a_x)-\tau_{yx}^{-}(-a_x)
	\notag\\
	&=\tau_{xy}^{+}(-a_y)-\tau_{xy}^{+}(a_y)+\tau_{xy}^{-}(a_y)-\tau_{xy}^{-}(-a_y),
	\label{YuzvyakEq17}
\end{align}
which forces the corner relations
\begin{equation}\label{YuzvyakEq18}
	\begin{array}{ll}
	\tau_{xy}^{+}(a_y)=\tau_{yx}^{+}(a_x)=\tau_{+}^{+},&\quad
	\tau_{xy}^{+}(-a_y)=\tau_{yx}^{-}(a_x)=\tau_{-}^{+},\\
	\tau_{xy}^{-}(a_y)=\tau_{yx}^{+}(-a_x)=\tau_{+}^{-},&\quad
	\tau_{xy}^{-}(-a_y)=\tau_{yx}^{-}(-a_x)=\tau_{-}^{-},
	\end{array}	
\end{equation}
so that Eq.~\eqref{YuzvyakEq17} reduces to
\begin{equation}
	\iint_{\mathcal{D}}\mathcal{V}(x,y) dx dy=
	\tau_{+}^{-}-\tau_{+}^{+}+\tau_{-}^{+}-\tau_{-}^{-}.
	\label{YuzvyakEq19}
\end{equation}

%-----------------------------------------------------------------
%--- Explicit shear stress
A second integration of Eq.~\eqref{YuzvyakEq13} with respect to \(y\) and
Eq.~\eqref{YuzvyakEq15} with respect to \(x\) provides the shear stress:
\begin{align}
	\sigma_{xy}(x,y)&=
	\frac{\tau_{xy}^{+}(y)+\tau_{xy}^{-}(y)+\tau_{yx}^{+}(x)+\tau_{yx}^{-}(x)}{2}
	-\frac{\tau_{+}^{+}+\tau_{+}^{-}+\tau_{-}^{+}+\tau_{-}^{-}}{4}
	\notag\\
	&\qquad-\frac{1}{4}\iint_{\mathcal{D}}
	\mathcal{V}(\xi,\eta) 
	\operatorname{sgn}(x-\xi)\operatorname{sgn}(y-\eta) d\xi d\eta .
	\label{YuzvyakEq20}
\end{align}

Formula \eqref{YuzvyakEq20} expresses the shear stress through the
Vihak function.  It satisfies the traction conditions
$(\ref{YuzvyakEq2})_{2,4}$ whenever the sectional resultants
\eqref{YuzvyakEq14}, \eqref{YuzvyakEq16}, the corner couplings
\eqref{YuzvyakEq18}, and the domain‐resultant condition
\eqref{YuzvyakEq19} are fulfilled.


Solving the remaining Poisson problems
\eqref{YuzvyakEq11}–\eqref{YuzvyakEq12} in the same manner gives the
normal stresses
\begin{align}
	\sigma_{xx}(x,y)
	&=-\frac{p_{x}^{+}(y)+p_{x}^{-}(y)}{2}
	-\frac{x-a_{x}}{2} \frac{d\tau_{xy}^{+}(y)}{dy}
	-\frac{x+a_{x}}{2} \frac{d\tau_{xy}^{-}(y)}{dy} \notag\\
	&\qquad
	+\frac{1}{2}\int_{-a_{x}}^{a_{x}}
	\mathcal{V}(\xi,y) \lvert x-\xi\rvert d\xi,
	\label{YuzvyakEq21}\\[3pt]
	%
	\sigma_{yy}(x,y)
	&=-\frac{p_{y}^{+}(x)+p_{y}^{-}(x)}{2}
	-\frac{y-a_{y}}{2} \frac{d\tau_{yx}^{+}(x)}{dx}
	-\frac{y+a_{y}}{2} \frac{d\tau_{yx}^{-}(x)}{dx} \notag\\
	&\qquad
	+\frac{1}{2}\int_{-a_{y}}^{a_{y}}
	\mathcal{V}(x,\eta) \lvert y-\eta\rvert d\eta.
	\label{YuzvyakEq22}
\end{align}

\noindent which comply with the mixed boundary data
$(\ref{YuzvyakEq2})_1,(\ref{YuzvyakEq8})_1$ and
$(\ref{YuzvyakEq2})_3,(\ref{YuzvyakEq8})_2$ provided
\begin{align}
	\int_{-a_{x}}^{a_{x}} x \mathcal{V}(x,y) dx&=
	p_{x}^{+}(y)-p_{x}^{-}(y)
	-a_{x} \frac{d}{dy}\bigl[\tau_{xy}^{+}(y)+\tau_{xy}^{-}(y)\bigr],
	\label{YuzvyakEq23}
\\
	\int_{-a_{y}}^{a_{y}} y \mathcal{V}(x,y) dy&=
	p_{y}^{+}(x)-p_{y}^{-}(x)
	-a_{y} \frac{d}{dx}\bigl[\tau_{yx}^{+}(x)+\tau_{yx}^{-}(x)\bigr].
	\label{YuzvyakEq24}
\end{align}
Relations \eqref{YuzvyakEq23}–\eqref{YuzvyakEq24} represent the first
moments of the Vihak function in the horizontal and vertical
cross-sections; higher-order moments are discussed next.

The sectional constraints \eqref{YuzvyakEq14} and \eqref{YuzvyakEq16}
give the \emph{resultants} of the Vihak function in the
$y = \text{const}$ and $x = \text{const}$ cross-sections,
respectively.  
The relations \eqref{YuzvyakEq23} and \eqref{YuzvyakEq24} provide the
corresponding first \emph{moments}, whereas
\eqref{YuzvyakEq19} is the domain-wide resultant.
Combining Eqs.~\eqref{YuzvyakEq23}–\eqref{YuzvyakEq24} with the corner
couplings \eqref{YuzvyakEq18} yields the higher-order identities

\begin{align}
	\iint_{\mathcal D} y \mathcal V dx dy
	&=\int_{-a_y}^{a_y} \bigl[\tau_{xy}^{+}(y)-\tau_{xy}^{-}(y)\bigr] dy
	-a_y\bigl(\tau_{+}^{+}+\tau_{-}^{+}-\tau_{+}^{-}-\tau_{-}^{-}\bigr)
	\notag\\
	&=\int_{-a_x}^{a_x} \bigl[p_{y}^{+}(x)-p_{y}^{-}(x)\bigr] dx
	-a_y\bigl(\tau_{+}^{+}+\tau_{-}^{+}-\tau_{+}^{-}-\tau_{-}^{-}\bigr),
	\label{YuzvyakEq24y}
	\\[6pt]
	%
	\iint_{\mathcal D} x \mathcal V dx dy
	&=\int_{-a_x}^{a_x} \bigl[\tau_{yx}^{+}(x)-\tau_{yx}^{-}(x)\bigr] dx
	-a_x\bigl(\tau_{+}^{+}+\tau_{+}^{-}-\tau_{-}^{+}-\tau_{-}^{-}\bigr)
	\notag\\
	&=\int_{-a_y}^{a_y} \bigl[p_{x}^{+}(y)-p_{x}^{-}(y)\bigr] dy
	-a_x\bigl(\tau_{+}^{+}+\tau_{+}^{-}-\tau_{-}^{+}-\tau_{-}^{-}\bigr),
	\label{YuzvyakEq24x}
	\\[6pt]
	%
	\iint_{\mathcal D} x y \mathcal V dx dy
	&=\int_{-a_y}^{a_y} y\bigl[p_{x}^{+}(y)-p_{x}^{-}(y)\bigr] dy
	+a_x \int_{-a_y}^{a_y} \bigl[\tau_{xy}^{+}(y)+\tau_{xy}^{-}(y)\bigr] dy
	\notag\\
	&\qquad\qquad-a_xa_y\bigl(\tau_{+}^{+}+\tau_{-}^{+}+\tau_{+}^{-}+\tau_{-}^{-}\bigr)\notag\\
	&=\int_{-a_x}^{a_x} x\bigl[p_{y}^{+}(x)-p_{y}^{-}(x)\bigr] dx
	+a_y \int_{-a_x}^{a_x} \bigl[\tau_{yx}^{+}(x)+\tau_{yx}^{-}(x)\bigr] dx\notag\\
	&\qquad\qquad-a_ya_x\bigl(\tau_{+}^{+}+\tau_{-}^{+}+\tau_{+}^{-}+\tau_{-}^{-}\bigr).
	\label{YuzvyakEq24xy}
\end{align}

From Eqs.~\eqref{YuzvyakEq24y}–\eqref{YuzvyakEq24xy} we deduce three
global-equilibrium conditions for the boundary tractions:

\begin{align}
	&\int_{-a_x}^{a_x} \bigl[p_{y}^{+}(x)-p_{y}^{-}(x)\bigr] dx
	=\int_{-a_y}^{a_y} \bigl[\tau_{xy}^{+}(y)-\tau_{xy}^{-}(y)\bigr] dy,
	\label{YuzvyakEq25}
	\\
	&\int_{-a_y}^{a_y} \bigl[p_{x}^{+}(y)-p_{x}^{-}(y)\bigr] dy
	=\int_{-a_x}^{a_x} \bigl[\tau_{yx}^{+}(x)-\tau_{yx}^{-}(x)\bigr] dx,
	\label{YuzvyakEq26}
	\\
	&\int_{-a_y}^{a_y} y\bigl[p_{x}^{+}(y)-p_{x}^{-}(y)\bigr] dy
	+a_x \int_{-a_y}^{a_y} \bigl[\tau_{xy}^{+}(y)+\tau_{xy}^{-}(y)\bigr] dy \notag
	\\
	&\qquad=\int_{-a_x}^{a_x} x\bigl[p_{y}^{+}(x)-p_{y}^{-}(x)\bigr] dx
	+a_y \int_{-a_x}^{a_x} \bigl[\tau_{yx}^{+}(x)+\tau_{yx}^{-}(x)\bigr] dx.
	\label{YuzvyakEq27}
\end{align}


The stress representations \eqref{YuzvyakEq20}–\eqref{YuzvyakEq22} now
depend solely on the Vihak potential
\(\mathcal{V}\) \eqref{YuzvyakEq9}.  Inserting them into the
compatibility relation \eqref{YuzvyakEq7} turns the original mixed
boundary-value problem into a single governing equation for
\(\mathcal{V}\), subject to the integral constraints
\eqref{YuzvyakEq14}, \eqref{YuzvyakEq16},
\eqref{YuzvyakEq23}, and \eqref{YuzvyakEq24}.  Consistency of those
constraints is guaranteed by the corner couplings
\eqref{YuzvyakEq18} together with the global-equilibrium conditions
\eqref{YuzvyakEq25}–\eqref{YuzvyakEq27}.  The same integral
relations were derived earlier from overall-equilibrium considerations
for a rectangular plate \cite{Yuzvyak32,Yuzvyak33}.

\section{Solution Representation}

To solve the boundary–value problem for the Vihak function we apply
variable separation with the complete systems of eigen- and associated
functions  
\(\{1,t,\cos(\gamma_{t,n}t),\sin(\lambda_{t,n}t)\}\)  
(\(t=x,y;\;n=1,2,\dots\)) introduced in
\cite{Yuzvyak29,Yuzvyak30}.  Here
\(\gamma_{t,n}=n\pi/a_{t}\) and \(\lambda_{t,n}>0\) is the \(n\)-th root
of \(\tan(\lambda_{t}a_{t})=\lambda_{t}a_{t}\).
For a generic field \(f(x,y)\) the expansion in the \(x\)-direction
reads
\begin{equation}
	f(x,y)=f^{0}(x,y)+f^{\mathrm e}(x,y),
	\label{YuzvyakEq28}
\end{equation}
with the \emph{elementary} part
\begin{equation}
	f^{0}(x,y)=
	\frac{1}{4a_{x}}\int_{-a_{x}}^{a_{x}}f(x,y) dx
	+\frac{3x}{4a_{x}}\int_{-a_{x}}^{a_{x}}x f(x,y) dx,
	\label{YuzvyakEq29}
\end{equation}
and the \emph{self-equilibrated} part
\begin{equation}
	f^{\mathrm e}(x,y)=
	\sum_{n=1}^{\infty} \Bigl[
	f_{n}^{\mathrm c}(y)\cos(\gamma_{x,n}x)
	+f_{n}^{\mathrm s}(y)\sin(\lambda_{x,n}x)
	\Bigr],
	\label{YuzvyakEq30}
\end{equation}
where
\begin{equation}\label{YuzvyakEq31}
\begin{array}{ll}
	f_{n}^{\mathrm c}(y)&=\frac{1}{a_{x}}
	\int_{-a_{x}}^{a_{x}}f(x,y)\cos(\gamma_{x,n}x) dx,
	\\
	f_{n}^{\mathrm s}(y)&=\frac{1}{a_{x}\sin^{2}(\lambda_{x,n}a_{x})}
	\int_{-a_{x}}^{a_{x}}f(x,y)\sin(\lambda_{x,n}x) dx.
\end{array}
\end{equation}

\emph{Decomposition of boundary tractions.}
Using Eqs.~\eqref{YuzvyakEq28}–\eqref{YuzvyakEq31} we split the prescribed
tractions into elementary and self-equilibrated parts:
\begin{align}
	p_{x}^{\pm}(y)&=p_{x}^{\pm0}(y)+p_{x}^{\pm\mathrm e}(y), &
	p_{y}^{\pm}(x)&=p_{y}^{\pm0}(x)+p_{y}^{\pm\mathrm e}(x),\notag\\
	\tau_{xy}^{\pm}(y)&=\tau_{xy}^{\pm0}(y)+\tau_{xy}^{\pm\mathrm e}(y), &
	\tau_{yx}^{\pm}(x)&=\tau_{yx}^{\pm0}(x)+\tau_{yx}^{\pm\mathrm e}(x).
	\label{YuzvyakEq32}
\end{align}

\noindent
Elementary parts are
\begin{equation}
	\begin{array}{ll}
		p_x^{\pm0}(y)
		&= \frac{1}{2a_y}\int_{-a_y}^{a_y}p_x^\pm(y) dy
		+ \frac{3y}{2a_y^3}\int_{-a_y}^{a_y}y p_x^\pm(y) dy,
		\\[8pt]
		p_y^{\pm0}(x)
		&= \frac{1}{2a_x}\int_{-a_x}^{a_x}p_y^\pm(x) dx
		+ \frac{3x}{2a_x^3}\int_{-a_x}^{a_x}x p_y^\pm(x) dx,
		\\[8pt]
		\tau_{xy}^{\pm0}(y)
		&= \frac{1}{2}\big[\tau_{xy}^{\pm}(a_y)+\tau_{xy}^{\pm}(-a_y)\big]
		+ \frac{y}{2a_y}\bigl[\tau_{xy}^{\pm}(a_y)-\tau_{xy}^{\pm}(-a_y)\bigr]
		\\
		&\quad+ \frac{3 (y^2-a_y^2)}{4a_y^2}\left[
		\tau_{xy}^{\pm}(a_y)+\tau_{xy}^{\pm}(-a_y)
		-\frac{1}{a_y}\int_{-a_y}^{a_y}\tau_{xy}^{\pm}(y) dy
		\right],
		\\[8pt]
		\tau_{yx}^{\pm0}(x)
		&= \frac{1}{2}\big[\tau_{yx}^{\pm}(a_x)+\tau_{yx}^{\pm}(-a_x)\big]
		 + \frac{x}{2a_x}\bigl[\tau_{yx}^{\pm}(a_x)-\tau_{yx}^{\pm}(-a_x)\bigr]
		\\
		&\quad+ \frac{3 (x^2-a_x^2)}{4a_x^2}\left[
		\tau_{yx}^{\pm}(a_x)+\tau_{yx}^{\pm}(-a_x)
		-\frac{1}{a_x}\int_{-a_x}^{a_x}\tau_{yx}^{\pm}(x) dx
		\right].
	\end{array}
	\label{YuzvyakEq32a}
\end{equation}

\noindent
Self-equilibrated parts read
\begin{equation}\label{YuzvyakEq32b}
	\begin{aligned}
		p_{x}^{\pm\mathrm e}(y)&=
		{\textstyle\sum_{m=1}^{\infty}}
		\bigl(p_{x\mathrm c,m}^{\pm}\cos\gamma_{y,m}y
		+p_{x\mathrm s,m}^{\pm}\sin\lambda_{y,m}y\bigr),
		\\
		p_{y}^{\pm\mathrm e}(x)&=
		{\textstyle\sum_{n=1}^{\infty}}
		\bigl(p_{y\mathrm c,n}^{\pm}\cos\gamma_{x,n}x
		+p_{y\mathrm s,n}^{\pm}\sin\lambda_{x,n}x\bigr),
		\\
		\tau_{xy}^{\pm\mathrm e}(y)&=
		{\textstyle\sum_{m=1}^{\infty}}
		\Bigl(
		\tau_{xy\mathrm c,m}^{\pm}\tfrac{\sin\gamma_{y,m}y}{\gamma_{y,m}}
		+\tau_{xy\mathrm s,m}^{\pm}
		\tfrac{\cos\lambda_{y,m}a_{y}-\cos\lambda_{y,m}y}{\lambda_{y,m}}
		\Bigr),
		\\
		\tau_{yx}^{\pm\mathrm e}(x)&=
		{\textstyle\sum_{n=1}^{\infty}}
		\Bigl(
		\tau_{yx\mathrm c,n}^{\pm}\tfrac{\sin\gamma_{x,n}x}{\gamma_{x,n}}
		+\tau_{yx\mathrm s,n}^{\pm}
		\tfrac{\cos\lambda_{x,n}a_{x}-\cos\lambda_{x,n}x}{\lambda_{x,n}}
		\Bigr).
	\end{aligned}
\end{equation}

Fourier coefficients (e.g. \(p_{x\mathrm c,m}^{\pm}\),  
\(\tau_{xy\mathrm c,m}^{\pm}\)) follow directly from
Eq.~\eqref{YuzvyakEq31}; for instance
\[
\begin{array}{ll}
	\tau_{xy\mathrm c ,m}^{\pm}
	&=\tfrac{\gamma_{y,m}}{a_y}
	\int_{-a_y}^{a_y} \tau_{xy}^{\pm}(y)\sin(\gamma_{y,m}y) dy,\\[3pt]
	\tau_{xy\mathrm s ,m}^{\pm}
	&=\tfrac{2\lambda_{y,m}^{2}}
	{ 2\lambda_{y,m}a_y-\sin(2\lambda_{y,m}a_y) }
	\int_{-a_y}^{a_y} \tau_{xy}^{\pm}(y)
	\bigl[\cos(\lambda_{y,m}a_y)-\cos(\lambda_{y,m}y)\bigr] dy,\\[3pt]
	\tau_{yx\mathrm c ,n}^{\pm}
	&=\tfrac{\gamma_{x,n}}{a_x}
	\int_{-a_x}^{a_x} \tau_{yx}^{\pm}(x)\sin(\gamma_{x,n}x) dx,\\[3pt]
	\tau_{yx\mathrm s ,n}^{\pm}
	&=\tfrac{2\lambda_{x,n}^{2}}
	{ 2\lambda_{x,n}a_x-\sin(2\lambda_{x,n}a_x) }
	\int_{-a_x}^{a_x} \tau_{yx}^{\pm}(x)
	\bigl[\cos(\lambda_{x,n}a_x)-\cos(\lambda_{x,n}x)\bigr] dx.
\end{array}
\]

%-----------------------------------------------------------------
\emph{Decomposition of additional strains.}
The additional-strain set
\[\Phi=\left\{\varepsilon_{xx}^{0}, \varepsilon_{yy}^{0},
\partial^{2}\varepsilon_{xy}^{0}/\partial x\partial y\right\}\]
is expanded as
\begin{align}
	\Phi={}&
	\Phi_{00}+x\Phi_{10}+y\Phi_{01}+xy\Phi_{11}\notag\\
	&+\sum_{n=1}^{\infty}
	\bigl[(\Phi_{n}^{\mathrm c0}+y\Phi_{n}^{\mathrm c1})\cos\gamma_{x,n}x
	+(\Phi_{n}^{\mathrm s0}+y\Phi_{n}^{\mathrm s1})\sin\lambda_{x,n}x
	\bigr]\notag\\
	&+\sum_{m=1}^{\infty}
	\bigl[(\Phi_{m}^{0\mathrm c}+x\Phi_{m}^{1\mathrm c})\cos\gamma_{y,m}y
	+(\Phi_{m}^{0\mathrm s}+x\Phi_{m}^{1\mathrm s})\sin\lambda_{y,m}y
	\bigr]\notag\\
	&+\sum_{n=1}^{\infty}\sum_{m=1}^{\infty}
	\Bigl[
	\Phi_{nm}^{\mathrm{cc}}\cos\gamma_{x,n}x \cos\gamma_{y,m}y
	+\Phi_{nm}^{\mathrm{cs}}\cos\gamma_{x,n}x \sin\lambda_{y,m}y\notag\\
	&\qquad+\Phi_{nm}^{\mathrm{sc}}\sin\lambda_{x,n}x \cos\gamma_{y,m}y
	+\Phi_{nm}^{\mathrm{ss}}\sin\lambda_{x,n}x \sin\lambda_{y,m}y
	\Bigr].
	\label{YuzvyakEq33}
\end{align}

%-----------------------------------------------------------------
\emph{Partition of the solution.}
With \eqref{YuzvyakEq33} we split the Vihak function and stresses into
three mutually orthogonal groups,
\begin{align}
	\mathcal V(x,y)&=\mathcal V^{0}(x,y)+\mathcal V^{\mathrm e}(x,y)
	+\tilde{\mathcal V}(x,y),
	\label{YuzvyakEq34}\\
	\sigma_{ij}(x,y)&=\sigma_{ij}^{0}(x,y)+\sigma_{ij}^{\mathrm e}(x,y)
	+\tilde{\sigma}_{ij}(x,y)\qquad (i,j=x,y),
	\label{YuzvyakEq35}
\end{align}
where  

* \(^{0}\) : elementary terms generated by the associated functions
\(1,x,y\);  

* \(^{\mathrm e}\) : self-equilibrated terms coming from the single
Fourier sums;  

* \(\tilde{\phantom{V}}\) : hybrid terms arising from the double Fourier
sums.

Because the governing equations are linear, the three subsystems can be
solved independently and superposed to construct the complete solution
for the rectangular domain~\(\mathcal D\).

\subsection{Elementary Solution}

First, we construct the elementary part of the Vihak function
\begin{equation}
	\mathcal V^{0}(x,y)
	= \mathcal V_{0}(x,y)
	+ x \mathcal V_{1}(x,y)
	+ y \mathcal V_{2}(x,y)
	+ xy \mathcal V_{3}(x,y),
	\label{YuzvyakEq36}
\end{equation}
corresponding to the decomposition in~Eq.~\eqref{YuzvyakEq34}.
Here
\begin{align*}
	\mathcal V_{0}(x,y)
	&= \tfrac{1}{4 a_{x} a_{y}}
	{\textstyle\iint_{\mathcal D}}  \mathcal V(x,y) dx dy
	= \tfrac{1}{4 a_{x} a_{y}}
	\bigl(\tau_{-}^{+}-\tau_{+}^{+}+\tau_{+}^{-}-\tau_{-}^{-}\bigr),
	\\[6pt]
	%
	\mathcal V_{1}(x,y)
	&= \tfrac{3}{4 a_{x}^{3} a_{y}}
	{\textstyle\iint_{\mathcal D}}  x \mathcal V(x,y) dx dy
	\\
	&= \tfrac{3}{4 a_{x}^{3} a_{y}}
	\Bigl[
	{\textstyle\int_{-a_{y}}^{a_{y}}} \bigl(p_{x}^{+}(y)-p_{x}^{-}(y)\bigr) dy
	- a_{x}\bigl(\tau_{+}^{+}+\tau_{+}^{-}-\tau_{-}^{+}-\tau_{-}^{-}\bigr)
	\Bigr]
	\\
	&= \tfrac{3}{4 a_{x}^{3} a_{y}}
	\Bigl[
	{\textstyle\int_{-a_{x}}^{a_{x}}} \bigl(\tau_{yx}^{+}(x)-\tau_{yx}^{-}(x)\bigr) dx
	- a_{x}\bigl(\tau_{+}^{+}+\tau_{+}^{-}-\tau_{-}^{+}-\tau_{-}^{-}\bigr)
	\Bigr],
	\\[6pt]
	%
	\mathcal V_{2}(x,y)
	&= \tfrac{3}{4 a_{y}^{3} a_{x}}
	{\textstyle\iint_{\mathcal D}}  y \mathcal V(x,y) dx dy
	\\
	&= \tfrac{3}{4 a_{y}^{3} a_{x}}
	\Bigl[
	{\textstyle\int_{-a_{x}}^{a_{x}}} \bigl(p_{y}^{+}(x)-p_{y}^{-}(x)\bigr) dx
	- a_{y}\bigl(\tau_{+}^{+}-\tau_{+}^{-}+\tau_{-}^{+}-\tau_{-}^{-}\bigr)
	\Bigr]
	\\
	&= \tfrac{3}{4 a_{y}^{3} a_{x}}
	\Bigl[
	{\textstyle\int_{-a_{y}}^{a_{y}}} \bigl(\tau_{xy}^{+}(y)-\tau_{xy}^{-}(y)\bigr) dy
	- a_{y}\bigl(\tau_{+}^{+}+\tau_{-}^{+}-\tau_{+}^{-}-\tau_{-}^{-}\bigr)
	\Bigr],
	\\[6pt]
	%
	\mathcal V_{3}(x,y)
	&= \tfrac{9}{4 a_{x}^{3} a_{y}^{3}}
	{\textstyle\iint_{\mathcal D}}  x y \mathcal V(x,y) dx dy
	\\
	&= \tfrac{9}{4 a_{x}^{3} a_{y}^{3}}
	\Bigl[
	{\textstyle\int_{-a_{y}}^{a_{y}}}  y\bigl(p_{x}^{+}(y)-p_{x}^{-}(y)\bigr) dy
	+ a_{x} {\textstyle\int_{-a_{y}}^{a_{y}}} \bigl(\tau_{xy}^{+}(y)+\tau_{xy}^{-}(y)\bigr) dy
	\\[-2pt]
	&\qquad
	- a_{x} a_{y}\bigl(\tau_{+}^{+}+\tau_{-}^{+}+\tau_{+}^{-}+\tau_{-}^{-}\bigr)
	\Bigr]
	\\
	&= \tfrac{9}{4 a_{x}^{3} a_{y}^{3}}
	\Bigl[
	{\textstyle\int_{-a_{x}}^{a_{x}}}  x\bigl(p_{y}^{+}(x)-p_{y}^{-}(x)\bigr) dx
	+ a_{y} {\textstyle\int_{-a_{x}}^{a_{x}}} \bigl(\tau_{yx}^{+}(x)+\tau_{yx}^{-}(x)\bigr) dx
	\\
	&\qquad
	- a_{x} a_{y}\bigl(\tau_{+}^{+}+\tau_{-}^{+}+\tau_{+}^{-}+\tau_{-}^{-}\bigr)
	\Bigr].
\end{align*}

\noindent These expressions follow from the decomposition algorithm
\eqref{YuzvyakEq28}--\eqref{YuzvyakEq31}
and the integral conditions
\eqref{YuzvyakEq17}, \eqref{YuzvyakEq24y}--\eqref{YuzvyakEq24xy}.

By substituting Eq.~\eqref{YuzvyakEq36} into the stress representations
\eqref{YuzvyakEq20}--\eqref{YuzvyakEq22}, and considering the ``0'' parts of
Eq.~\eqref{YuzvyakEq32}, the compatibility equation \eqref{YuzvyakEq7} implies
$\mathcal V_{\ell}=0$, $\ell=0,1,2,3$, which leads to the conditions
\begin{align}
	a_{x}\tau_{y}
	&= a_{x}\bigl(\tau_{+}^{+}+\tau_{+}^{-}-\tau_{-}^{+}-\tau_{-}^{-}\bigr)\notag\\
	&= {\textstyle\int_{-a_{x}}^{a_{x}}} (\tau_{yx}^{+}-\tau_{yx}^{-}) dx
	\notag\\
	&= {\textstyle\int_{-a_{y}}^{a_{y}}} (p_{x}^{+}-p_{x}^{-}) dy,
	\notag\\[4pt]
	%
	a_{y}\tau_{x}
	&= a_{y}\bigl(\tau_{+}^{+}+\tau_{-}^{+}-\tau_{+}^{-}-\tau_{-}^{-}\bigr)\notag\\
	&= {\textstyle\int_{-a_{y}}^{a_{y}}} (\tau_{xy}^{+}-\tau_{xy}^{-}) dy\notag\\
	&= {\textstyle\int_{-a_{x}}^{a_{x}}} (p_{y}^{+}-p_{y}^{-}) dx,
	\notag\\[4pt]
	%
	a_{x}a_{y}\tau
	&= a_{x}a_{y}\bigl(\tau_{+}^{+}+\tau_{-}^{+}+\tau_{+}^{-}+\tau_{-}^{-}\bigr)\notag\\
	&= {\textstyle\int_{-a_{y}}^{a_{y}}}  y (p_{x}^{+}-p_{x}^{-}) dy
	+ a_{x} {\textstyle\int_{-a_{y}}^{a_{y}}} (\tau_{xy}^{+}+\tau_{xy}^{-}) dy\notag\\
	&= {\textstyle\int_{-a_{x}}^{a_{x}}}  x (p_{y}^{+}-p_{y}^{-}) dx
	+ a_{y} {\textstyle\int_{-a_{x}}^{a_{x}}} (\tau_{yx}^{+}+\tau_{yx}^{-}) dx,
	\notag\\[4pt]
	%
	\tau_{+}^{+}-\tau_{-}^{+}-\tau_{+}^{-}+\tau_{-}^{-} &= 0.
	\label{YuzvyakEq37}
\end{align}


These conditions ensure the \textit{compatibility of the applied loadings}.
Hence, $\mathcal V^{0}\equiv 0$, and the elementary stress components
\eqref{YuzvyakEq35} are
\begin{align}
	\sigma_{xx}^{0}(x,y)
	&= -\frac{a_{x}+x}{2a_{x}} p_{x}^{+0}(y)
	-\frac{a_{x}-x}{2a_{x}} p_{x}^{-0}(y),
	\label{YuzvyakEq38}
	\\
	\sigma_{yy}^{0}(x,y)
	&= -\frac{a_{y}+y}{2a_{y}} p_{y}^{+0}(x)
	-\frac{a_{y}-y}{2a_{y}} p_{y}^{-0}(x),
	\label{YuzvyakEq39}
	\\
	\sigma_{xy}^{0}(x,y)
	&= \frac{1}{4}\left(\tau + x \frac{\tau_{x}}{a_{x}}
	+ y \frac{\tau_{y}}{a_{y}}\right)
	\notag\\
	&\qquad
	+ \frac{3(x^{2}-a_{x}^{2})}{8 a_{x}^{3}}
	\left[a_{x}\tau-\int_{-a_{x}}^{a_{x}} \bigl(\tau_{yx}^{+}+\tau_{yx}^{-}\bigr) dx\right]\nonumber
	\\
	&\qquad
	+ \frac{3(y^{2}-a_{y}^{2})}{8 a_{y}^{3}}
	\left[a_{y}\tau-\int_{-a_{y}}^{a_{y}} \bigl(\tau_{xy}^{+}+\tau_{xy}^{-}\bigr) dy\right].
	\label{YuzvyakEq40}
\end{align}
These results coincide with those reported in
\cite{Yuzvyak30,Yuzvyak31,Yuzvyak32}
and satisfy the global equilibrium conditions.

\subsection{Self-equilibrated Solution} 

For constructing the self-equilibrated parts of the solution
\eqref{YuzvyakEq34} and \eqref{YuzvyakEq35}, 
we involve only the self-equilibrated parts of the external loadings, 
which satisfy the following equilibrium conditions:
\begin{align}
	\int_{-a_\xi}^{a_\xi} p_\eta^{\pm\mathrm e}(\xi) d\xi
	&= \int_{-a_\xi}^{a_\xi} \xi p_\eta^{\pm\mathrm e}(\xi) d\xi
	= \int_{-a_\xi}^{a_\xi} \tau_{\eta\xi}^{\pm\mathrm e}(\xi) d\xi
	\notag\\
	&= \tau_{\eta\xi}^{\pm\mathrm e}(\pm a_\eta)
	= 0
	\quad (\eta,\xi = x,y,\ \eta\neq\xi).
	\label{YuzvyakEq41}
\end{align}

Based on the decomposition algorithm 
\eqref{YuzvyakEq28}--\eqref{YuzvyakEq31}, 
within the context of Eq.~\eqref{YuzvyakEq41}, 
the self-equilibrated part of the Vihak function is
\begin{align}
	\mathcal V^{\mathrm e}(x,y)
	&= -\frac{1}{2a_x}\frac{d}{dy}\Bigl(
	\tau_{xy}^{+\mathrm e}(y)-\tau_{xy}^{-\mathrm e}(y)\Bigr)
	-\frac{1}{2a_y}\frac{d}{dx}\Bigl(
	\tau_{yx}^{+\mathrm e}(x)-\tau_{yx}^{-\mathrm e}(x)\Bigr)\notag\\
	&\quad +\frac{3x}{2a_x^3}\Bigl[
	p_x^{+\mathrm e}(y)-p_x^{-\mathrm e}(y)
	-a_x\frac{d}{dy}\bigl(\tau_{xy}^{+\mathrm e}(y)+\tau_{xy}^{-\mathrm e}(y)\bigr)
	\Bigr]\notag\\
	&\quad +\frac{3y}{2a_y^3}\Bigl[
	p_y^{+\mathrm e}(x)-p_y^{-\mathrm e}(x)
	-a_y\frac{d}{dx}\bigl(\tau_{yx}^{+\mathrm e}(x)+\tau_{yx}^{-\mathrm e}(x)\bigr)
	\Bigr].
	\label{YuzvyakEq42}
\end{align}

The self-equilibrated parts of the stress-tensor components are then

\begin{align}
	\sigma_{xx}^{\mathrm e}(x,y)
	&= -\frac{(x+a_x)^2(2a_x-x)}{4a_x^3} p_x^{+\mathrm e}(y)
	-\frac{(x-a_x)^2(2a_x+x)}{4a_x^3} p_x^{-\mathrm e}(y) \notag\\
	&\quad -\frac{(x-a_x)^2}{4a_x^2}\biggl[
	(x+a_x)\frac{d\tau_{xy}^{+\mathrm e}}{dy}
	+(x-a_x)\frac{d\tau_{xy}^{-\mathrm e}}{dy}
	\biggr] \notag\\
	&\quad -\frac{1}{4a_y}
	\int_{-a_x}^{a_x} 
	\bigl(\tau_{yx}^{+\mathrm e}(\xi)-\tau_{yx}^{-\mathrm e}(\xi)\bigr)
	 \mathrm{sgn}(x-\xi) d\xi 
	\notag\\
	&\quad -\frac{3y}{4a_y^3}\biggl[
	a_y \int_{-a_x}^{a_x} 
	\bigl(\tau_{yx}^{+\mathrm e}(\xi)+\tau_{yx}^{-\mathrm e}(\xi)\bigr)
	\mathrm{sgn}(x-\xi) d\xi 
	\notag\\
	&\qquad\qquad
	+ \int_{-a_x}^{a_x} 
	\bigl(p_y^{+\mathrm e}(\xi)-p_y^{-\mathrm e}(\xi)\bigr)
	|x-\xi| d\xi
	\biggr], 
	\label{YuzvyakEq43}\\
	%
	\sigma_{yy}^{\mathrm e}(x,y)
	&= -\frac{(y+a_y)^2(2a_y-y)}{4a_y^3} p_y^{+\mathrm e}(x)
	-\frac{(y-a_y)^2(2a_y+y)}{4a_y^3} p_y^{-\mathrm e}(x) \notag\\
	&\quad -\frac{(y-a_y)^2}{4a_y^2}\biggl[
	(y+a_y)\frac{d\tau_{yx}^{+\mathrm e}}{dx}
	+(y-a_y)\frac{d\tau_{yx}^{-\mathrm e}}{dx}
	\biggr] \notag\\
	&\quad -\frac{1}{4a_x}
	\int_{-a_y}^{a_y} 
	\bigl(\tau_{xy}^{+\mathrm e}(\eta)-\tau_{xy}^{-\mathrm e}(\eta)\bigr)
	 \mathrm{sgn}(y-\eta) d\eta \notag\\
	&\quad -\frac{3x}{4a_x^3}\biggl[
	a_x \int_{-a_y}^{a_y} 
	\bigl(\tau_{xy}^{+\mathrm e}(\eta)+\tau_{xy}^{-\mathrm e}(\eta)\bigr)
	\mathrm{sgn}(y-\eta) d\eta \notag\\
	&\qquad\qquad
	- \int_{-a_y}^{a_y} 
	\bigl(p_x^{+\mathrm e}(\eta)-p_x^{-\mathrm e}(\eta)\bigr)
	|y-\eta| d\eta
	\biggr], 
	\label{YuzvyakEq44}\\
	%
	\sigma_{xy}^{\mathrm e}(x,y)
	&= \frac{a_y+y}{2a_y} \tau_{yx}^{+\mathrm e}(x)
	+\frac{a_y-y}{2a_y} \tau_{yx}^{-\mathrm e}(x)
	\notag\\
	&\quad+\frac{a_x+x}{2a_x} \tau_{xy}^{+\mathrm e}(y)
	+\frac{a_x-x}{2a_x} \tau_{xy}^{-\mathrm e}(y) 
	\notag\\
	&\quad +\frac{3(x^2-a_x^2)}{4a_x^3}\biggl[
	a_x\bigl(\tau_{xy}^{+\mathrm e}(y)+\tau_{xy}^{-\mathrm e}(y)\bigr)
	\notag\\
	&\qquad\qquad-\frac12\int_{-a_y}^{a_y} 
	\bigl(p_x^{+\mathrm e}(\eta)-p_x^{-\mathrm e}(\eta)\bigr)
	\mathrm{sgn}(y-\eta) d\eta
	\biggr] 
	\notag\\
	&\quad +\frac{3(y^2-a_y^2)}{4a_y^3}\biggl[
	a_y\bigl(\tau_{yx}^{+\mathrm e}(x)+\tau_{yx}^{-\mathrm e}(x)\bigr)
	\notag\\
	&\qquad\qquad-\frac12\int_{-a_x}^{a_x} 
	\bigl(p_y^{+\mathrm e}(\xi)-p_y^{-\mathrm e}(\xi)\bigr)
	\mathrm{sgn}(x-\xi) d\xi
	\biggr].
	\label{YuzvyakEq45}
\end{align}

The obtained expressions for the self‑equilibrated parts \eqref{YuzvyakEq43}--\eqref{YuzvyakEq45} of the stress‑tensor components \eqref{YuzvyakEq35} represent the local response of the structure to the self‑equilibrated external loadings given in Eq.~\eqref{YuzvyakEq32b}.

\subsection{Hybrid Solution}

The hybrid constituent $\tilde{\mathcal{V}}$ of Eq.~\eqref{YuzvyakEq34} satisfies the homogeneous integral conditions in both directions  
\begin{align}
	\int_{-a_x}^{a_x}\tilde{\mathcal{V}} dx
	=\int_{-a_x}^{a_x}x \tilde{\mathcal{V}} dx
	=\int_{-a_y}^{a_y}\tilde{\mathcal{V}} dy
	=\int_{-a_y}^{a_y}y \tilde{\mathcal{V}} dy
	=0,
	\label{YuzvyakEq46}
\end{align}
which follow from conditions \eqref{YuzvyakEq14}, \eqref{YuzvyakEq16}, \eqref{YuzvyakEq23}, and \eqref{YuzvyakEq24}.

Using Eqs.~\eqref{YuzvyakEq20}--\eqref{YuzvyakEq22}, the corresponding stress components in Eq.~\eqref{YuzvyakEq35} are
\begin{align}
	\tilde{\sigma}_{xx}(x,y)
	&=\frac{1}{2}\int_{-a_x}^{a_x}\tilde{\mathcal{V}}(\xi ,y) |x-\xi| d\xi,
	\notag\\
	\tilde{\sigma}_{yy}(x,y)
	&=\frac{1}{2}\int_{-a_y}^{a_y}\tilde{\mathcal{V}}(x,\eta) |y-\eta| d\eta,
	\notag\\
	\tilde{\sigma}_{xy}(x,y)
	&=-\frac{1}{4}\iint_{\mathcal D}
	\tilde{\mathcal{V}}(\xi ,\eta )
	 \mathrm{sgn}(x-\xi) \mathrm{sgn}(y-\eta)
	 d\xi d\eta.
	\label{YuzvyakEq47}
\end{align}

By substituting the decompositions \eqref{YuzvyakEq33} and \eqref{YuzvyakEq34} into the compatibility equation \eqref{YuzvyakEq7} within the context of the self‑equilibrated solution \eqref{YuzvyakEq42}, we obtain
\begin{align}
	&2\tilde{\mathcal{V}}(x,y)
	+\frac{1}{2}\int_{-a_x}^{a_x}
	\frac{\partial^2\tilde{\mathcal{V}}(\xi ,y)}{\partial y^2}
	 |x-\xi| d\xi
	+\frac{1}{2}\int_{-a_y}^{a_y}
	\frac{\partial^2\tilde{\mathcal{V}}(x,\eta)}{\partial x^2}
	 |y-\eta| d\eta
	\notag\\
	&\qquad =-E_0\Biggl(
	\frac{\partial^2\varepsilon_{xx}^0}{\partial y^2}
	+\frac{\partial^2\varepsilon_{yy}^0}{\partial x^2}
	-2\frac{\partial^2\varepsilon_{xy}^0}{\partial x \partial y}
	\Biggr)
	\notag\\
	&\qquad\quad-\frac{3x}{a_x^3}\bigl(p_x^{+\mathrm e}(y)-p_x^{-\mathrm e}(y)\bigr)
	+\frac{3x+a_x}{a_x^2}\frac{d\tau_{xy}^{+\mathrm e}(y)}{dy}
	+\frac{3x-a_x}{a_x^2}\frac{d\tau_{xy}^{-\mathrm e}(y)}{dy}
	\notag\\
	&\qquad\quad+\frac{d^2}{dy^2}\Biggl[
	\frac{(2a_x-x)(x+a_x)^2}{4a_x^3}p_x^{+\mathrm e}(y)
	+\frac{(2a_x+x)(x-a_x)^2}{4a_x^3}p_x^{-\mathrm e}(y)
	\notag\\
	&\qquad\quad
	+\frac{(x+a_x)(x-a_x)^2}{4a_x^2}\frac{d\tau_{xy}^{+\mathrm e}(y)}{dy}
	+\frac{(x-a_x)(x+a_x)^2}{4a_x^2}\frac{d\tau_{xy}^{-\mathrm e}(y)}{dy}
	\Biggr]
	\notag\\
	&\qquad\quad-\frac{3y}{a_y^3}\bigl(p_y^{+\mathrm e}(x)-p_y^{-\mathrm e}(x)\bigr)
	+\frac{3y+a_y}{a_y^2}\frac{d\tau_{yx}^{+\mathrm e}(x)}{dx}
	+\frac{3y-a_y}{a_y^2}\frac{d\tau_{yx}^{-\mathrm e}(x)}{dx}
	\notag\\
	&\qquad\quad+\frac{d^2}{dx^2}\Biggl[
	\frac{(2a_y-y)(y+a_y)^2}{4a_y^3}p_y^{+\mathrm e}(x)
	+\frac{(2a_y+y)(y-a_y)^2}{4a_y^3}p_y^{-\mathrm e}(x)
	\notag\\
	&\qquad\quad +\frac{(y+a_y)(y-a_y)^2}{4a_y^2}\frac{d\tau_{yx}^{+\mathrm e}(x)}{dx}
	+\frac{(y-a_y)(y+a_y)^2}{4a_y^2}\frac{d\tau_{yx}^{-\mathrm e}(x)}{dx}
	\Biggr],
	\label{YuzvyakEq48}
\end{align}
subject to the integral constraints \eqref{YuzvyakEq46}.

A solution to Eq.~\eqref{YuzvyakEq48} can be expressed in the closed form
\begin{align}
	\tilde{\mathcal{V}}(x,y)
	&=\sum_{n=1}^{\infty}\sum_{m=1}^{\infty}\Bigl[
	\mathcal{V}_{nm}^{cc}\cos\gamma_{x,n}x \cos\gamma_{y,m}y
	+\mathcal{V}_{nm}^{cs}\cos\gamma_{x,n}x \sin\lambda_{y,m}y
	\notag\\
	&\qquad
	+\mathcal{V}_{nm}^{sc}\sin\lambda_{x,n}x \cos\gamma_{y,m}y
	+\mathcal{V}_{nm}^{ss}\sin\lambda_{x,n}x \sin\lambda_{y,m}y
	\Bigr],
	\label{YuzvyakEq49}
\end{align}
and the corresponding stress‑tensor components follow by substitution into Eq.~\eqref{YuzvyakEq47}.

\section{Numerical Examples and Discussion}
\noindent
Consider the case study when the stress field in the rectangular domain $\mathcal{D}$ is induced solely by the thermal field, with the external boundary tractions set to zero.  
Assume the thermal strains to be
\begin{align}
	\varepsilon_T(x,y)=1-\frac{x^2y^2}{a_x^2a_y^2},
	\label{YuzvyakEq50}
\end{align}
and the material properties in the form
\begin{align}
	E=E_0\exp(\varkappa xy),\qquad
	\nu=\nu_0=0.3,\qquad
	E_0=\mathrm{const},
	\label{YuzvyakEq51}
\end{align}
where $\varkappa$ is the non‑homogeneity parameter.  
If $\varkappa=0$, the material is homogeneous; otherwise, the Young modulus varies in both planar coordinates. 

\begin{figure}[b]
	\centering
	\begin{subfigure}[t]{0.48\textwidth}
		\centering
		\includegraphics[width=\linewidth]{YuzvyakFig1a.png}
		\caption{Normal stress $\sigma_{xx}/E_0$ for the homogeneous material ($\varkappa=0$).}
		\label{YuzvyakFig1a}
	\end{subfigure}
	\hfill
	\begin{subfigure}[t]{0.48\textwidth}
		\centering
		\includegraphics[width=\linewidth]{YuzvyakFig1b.png}
		\caption{Tangential stress $\sigma_{xy}/E_0$ for the homogeneous material ($\varkappa=0$).}
		\label{YuzvyakFig1b}
	\end{subfigure}
	\caption{Full‑field distribution of the dimensionless stresses in the rectangular domain for the homogeneous material ($\varkappa=0$ in Eq.~\eqref{YuzvyakEq51}).}
	\label{YuzvyakFig1}
\end{figure}

\begin{figure}[t]
	\centering
	\begin{subfigure}[t]{0.48\textwidth}
		\centering
		\includegraphics[width=\linewidth]{YuzvyakFig2a.png}
		\caption{Normal stress $\sigma_{xx}/E_0$ for the inhomogeneous material ($\varkappa=1$).}
		\label{YuzvyakFig2a}
	\end{subfigure}
	\hfill
	\begin{subfigure}[t]{0.48\textwidth}
		\centering
		\includegraphics[width=\linewidth]{YuzvyakFig2b.png}
		\caption{Tangential stress $\sigma_{xy}/E_0$ for the inhomogeneous material ($\varkappa=1$).}
		\label{YuzvyakFig2b}
	\end{subfigure}
	\caption{Full‑field distribution of the dimensionless stresses in the rectangular domain for the inhomogeneous material ($\varkappa=1$ in Eq.~\eqref{YuzvyakEq51}).}
	\label{YuzvyakFig2}
\end{figure}


The normal and tangential stresses $\sigma_{xx}/E_0$ and $\sigma_{xy}/E_0$ computed in a square domain ($a_x=a_y=1$) are shown in Fig.~\ref{YuzvyakFig1} for a homogeneous material ($\varkappa=0$).  
Both stresses satisfy the boundary conditions: $\sigma_{xx}=0$ at $x=\pm a_x$ and $\sigma_{xy}=0$ on the entire periphery.  
The stress $\sigma_{xx}$ is nonzero at $y=\pm a_y$ due to thermal expansion and equilibrium requirements.  
Because of the symmetric thermal field, $\sigma_{yy}$ is obtained from $\sigma_{xx}$ by exchanging $x$ and $y$.  
The normal stresses are even in each coordinate, while $\sigma_{xy}$ is odd.

The effect of inhomogeneity ($\varkappa=1$) is illustrated in Fig.~\ref{YuzvyakFig2}.  
The stresses lose symmetry: $\sigma_{xx}$ is no longer even, and $\sigma_{xy}$ is no longer odd.  
Moreover, the peak of $\sigma_{xx}$ on the edges $y=\pm a_y$ increases, highlighting how inhomogeneity strongly affects the stress field under the same loading.

\medskip

Next, consider a rectangular domain $\mathcal{D}$ with $a_x=0.2$~m and $a_y=0.1$~m subjected to the force loadings
\begin{align}
	p_x^+(y) &= \frac{3p_0}{2}\left(1-\frac{y^2}{a_y^2}\right), 
	\quad p_x^-(y)=p_0,
	\notag\\
	p_y^\pm(x) &= \tau_{xy}^\pm(y) = \tau_{yx}^\pm(x) = 0,
	\label{YuzvyakEq52}
\end{align}
which satisfy the equilibrium conditions \eqref{YuzvyakEq25}--\eqref{YuzvyakEq27} for the resultant forces and moments.  

\begin{figure}[t]
	\centering
	\begin{subfigure}[t]{0.48\textwidth}
		\centering
		\includegraphics[width=\linewidth]{YuzvyakFig3a.png}
		\caption{Normal stress $\sigma_{yy}/E_0$ for a homogeneous material ($\nu_0=0.3$).}
		\label{YuzvyakFig3a}
	\end{subfigure}
	\hfill
	\begin{subfigure}[t]{0.48\textwidth}
		\centering
		\includegraphics[width=\linewidth]{YuzvyakFig3b.png}
		\caption{Tangential stress $\sigma_{xy}/E_0$ for a homogeneous material ($\nu_0=0.3$).}
		\label{YuzvyakFig3b}
	\end{subfigure}
	\caption{Full‑field distribution of the dimensionless stresses for a homogeneous material ($\nu_0=0.3$) under the force loading \eqref{YuzvyakEq52}.}
	\label{YuzvyakFig3}
\end{figure}


Because the side $y=-a_y$ carries a uniform pressure equal to the resultant of the non‑uniform load on $y=+a_y$, the stress distribution is more disturbed near the latter side (Fig.~\ref{YuzvyakFig3}).

\medskip

Now, consider an inhomogeneous material with the variable Young modulus
\begin{align}
	E = E_0\left(2 + 3 \frac{y^2}{a_y^2}\right),
	\label{YuzvyakEq53}
\end{align}
while $\nu_0=\mathrm{const}$.  
The corresponding stresses are shown in Fig.~\ref{YuzvyakFig4}.  
Compared to Fig.~\ref{YuzvyakFig3}, the variation of $E(y)$ causes a visible distortion of the stress field near $y=-a_y$ and alters the peak values of both $\sigma_{yy}$ and $\sigma_{xy}$.

\begin{figure}[t]
	\centering
	\begin{subfigure}[t]{0.48\textwidth}
		\centering
		\includegraphics[width=\linewidth]{YuzvyakFig4a.png}
		\caption{Normal stress $\sigma_{yy}/E_0$ for the inhomogeneous material.}
		\label{YuzvyakFig4a}
	\end{subfigure}
	\hfill
	\begin{subfigure}[t]{0.48\textwidth}
		\centering
		\includegraphics[width=\linewidth]{YuzvyakFig4b.png}
		\caption{Tangential stress $\sigma_{xy}/E_0$ for the inhomogeneous material.}
		\label{YuzvyakFig4b}
	\end{subfigure}
	\caption{Full‑field distribution of the dimensionless stresses for the inhomogeneous material with $E(y)$ from Eq.~\eqref{YuzvyakEq53} and $\nu_0=0.3$ under the force loading \eqref{YuzvyakEq52}.}
	\label{YuzvyakFig4}
\end{figure}


\medskip

Finally, consider a two‑layer plate $\mathcal D = \mathcal D_1 \cup \mathcal D_2$ with  
\[
\mathcal{D}_1=[-a_x,a_x]\times[-a_y,0],\quad
\mathcal{D}_2=[-a_x,a_x]\times[0,a_y],
\]
perfectly bonded along $y=0$.  
Let the piecewise material properties be represented as
\begin{align}
	f(x,y)= f_1(x,y) + \bigl(f_2(x,y)-f_1(x,y)\bigr)H(y),
	\label{YuzvyakEq54}
\end{align}
where $H(y)$ is the Heaviside step function and $f_\ell$ denotes either $E_\ell$ or $\nu_\ell$ in layer $\mathcal D_\ell$ ($\ell=1,2$).  
Assume
\[
\nu_1=0.3,\quad E_1=6.5E_0,\quad
\nu_2=0.28,\quad E_2=E_0\Bigl(2+11(y/a_y)^2\Bigr).
\]

\begin{figure}[t]
	\centering
	\begin{subfigure}[t]{0.48\textwidth}
		\centering
		\includegraphics[width=\linewidth]{YuzvyakFig5a.png}
		\caption{Normal stress $\sigma_{xx}/E_0$.}
		\label{YuzvyakFig5a}
	\end{subfigure}
	\hfill
	\begin{subfigure}[t]{0.48\textwidth}
		\centering
		\includegraphics[width=\linewidth]{YuzvyakFig5b.png}
		\caption{Normal stress $\sigma_{yy}/E_0$.}
		\label{YuzvyakFig5b}
	\end{subfigure}
	
	\begin{subfigure}[t]{0.48\textwidth}
		\centering
		\includegraphics[width=\linewidth]{YuzvyakFig5c.png}
		\caption{Tangential stress $\sigma_{xy}/E_0$.}
		\label{YuzvyakFig5c}
	\end{subfigure}
	\caption{Full‑field distribution of the dimensionless stresses in a two‑layer plate with material properties \eqref{YuzvyakEq54} under the loading \eqref{YuzvyakEq52}.}
	\label{YuzvyakFig5}
\end{figure}


Figure~\ref{YuzvyakFig5} shows the stress patterns for the two‑layer plate.  
Under the perfect bonding condition, $\sigma_{yy}$ and $\sigma_{xy}$ remain continuous across $y=0$, while $\sigma_{xx}$ exhibits a discontinuity at the interface.  
This discontinuity may trigger residual strains and, if unrelieved, potential fracture initiation.  
The most pronounced change in $\sigma_{yy}$ appears near $y=0$ and close to the vertical edges $x=\pm a_x$.

\section{Conclusions}
\noindent
A hybrid analytical–numerical algorithm has been developed for evaluating plane-stress/plane-strain fields in an arbitrarily inhomogeneous rectangular domain with corner points.  
The approach synergistically combines the \emph{method of additional (eigen)strains} with Vihak’s \emph{direct-integration method}.  

\begin{itemize}
	\item  The additional-strain formulation rewrites the constitutive law of an inhomogeneous solid in the form valid for a reference homogeneous medium, shifting all spatially varying elastic moduli into initially unknown strain fields.  
	\item  Applying direct integration then introduces a single scalar potential—the \emph{Vihak function}—whose governing boundary-value problem involves only a second-order integro-differential equation with integral side constraints.  This bypasses the usual difficulties of enforcing boundary conditions near the rectangle’s corners.  
	\item  Separation of variables in a complete system of eigen- and associated functions splits the solution into \emph{elementary}, \emph{self-equilibrated}, and \emph{hybrid} parts, each available in closed form.  The successive-approximation loop explicitly updates the additional strains, ensuring rapid convergence even for steep property gradients.  
\end{itemize}

Benchmark studies confirm the algorithm’s efficiency for (i) one- and two-dimensional material gradients and (ii) multilayer stacks with piecewise-constant or smoothly varying moduli.  Because a multilayer assembly is treated as a single domain with stepwise properties, the computational cost grows only mildly with the number of layers—making the method attractive for verifying homogenization or layered-composite models.

\bigskip
\section*{Acknowledgement}
This research was partially supported by the budget programme of Ukraine
``Support for the Development of Priority Research Areas''
(CPCEC 6451230) under grant No.~0125U000322
\emph{``Mathematical modelling and identification of protective surface and structural
	properties of materials with respect to mechanical, thermal, acoustic and electromagnetic influences''}.



\bibliographystyle{unsrtnat}
\bibliography{Chapter38}
